\chapter{Buccal Space Abscess}

\section*{Definition and Overview}
A buccal space abscess is a collection of pus within the buccal space, a potential tissue compartment of the cheek located between the buccinator muscle and the overlying skin. It almost always arises from an odontogenic (dental) infection, most commonly a decayed upper molar, that tracks from the tooth root into this space. The result is a progressively painful swelling of the cheek that can enlarge over days. It is one of several facial spaces that dental infections can involve, and it is important because untreated infection can spread to adjacent spaces, the orbit, or the neck. Prompt drainage and treatment of the source tooth are essential.

\section*{Epidemiology}
Buccal space abscesses are uncommon, and reliable incidence figures are lacking. They are most often seen in adults with untreated dental decay, particularly in the molar region, and can also occur in children with severe tooth infections. People with poor access to dental care, diabetes, or impaired immunity are at higher risk, as are those who delay seeking treatment for toothache. Because the condition reflects advanced dental disease, its frequency falls in populations with good preventive oral health care.

\section*{Aetiology and Pathophysiology}
The infection is dental in origin in the vast majority of cases. Bacteria from a decayed tooth reach the root apex and form a periapical abscess; when the pressure within the bone rises, the infection perforates the cortex and spreads into the surrounding facial spaces, in this case the buccal space.

\begin{itemize}
    \item \textbf{Odontogenic source:} an upper or lower molar with deep decay, pulp necrosis, or a fractured tooth
    \item \textbf{Direct inoculation:} less commonly, infection follows trauma to the cheek, a bite, an injection, or dental surgery
    \item \textbf{Microbiology:} mixed oral flora, including \textit{Streptococcus} and \textit{Staphylococcus} species together with anaerobes such as \textit{Fusobacterium} and \textit{Prevotella}
    \item \textbf{Mechanism:} the immune response produces an accumulation of dead tissue, bacteria, and white blood cells (pus) within the buccal space, causing the firm, painful swelling of the cheek
\end{itemize}

\section*{Clinical Features}
The dominant feature is a swelling of the cheek that typically develops over several days.

\begin{itemize}
    \item Fullness and swelling of the cheek, which may be firm initially and become fluctuant as pus accumulates
    \item Pain in the cheek that progressively worsens
    \item Toothache in a back tooth, often with visible decay or a fractured tooth
    \item Tenderness and redness of the skin overlying the swelling
    \item Trismus, or difficulty opening the mouth fully
    \item Fever, chills, and general malaise in some cases
    \item A history of recent toothache, dental treatment, or facial injury
\end{itemize}

\section*{Differential Diagnosis}
Other causes of cheek swelling should be considered before confirming the diagnosis.

\begin{itemize}
    \item \textbf{Other facial space infections:} submandibular, masticator, or canine-space abscesses
    \item \textbf{Parotitis:} infection of the parotid salivary gland
    \item \textbf{Jaw cysts:} an infected dentigerous or radicular cyst
    \item \textbf{Salivary gland stones:} causing obstructive swelling of a salivary gland
    \item \textbf{Trauma:} a haematoma within the cheek after injury
    \item \textbf{Tumours:} benign or malignant lesions of the cheek, jaw, or parotid gland
    \item \textbf{Lymph node infection:} suppurative lymphadenitis in the buccal region
\end{itemize}

\section*{When to Seek Medical Care}
Seek urgent dental or medical review for a painful tooth accompanied by swelling of the cheek, because dental infections can spread quickly. The sooner treatment begins, the simpler it is usually.

\textbf{Call 000 (or local emergency services) for:}
\begin{itemize}
    \item Swelling of the face or neck that is spreading rapidly
    \item Difficulty breathing or swallowing
    \item Inability to open the mouth fully
    \item High fever with rigors, confusion, or drowsiness
    \item Swelling spreading towards the eye, which can threaten vision
    \item Signs of sepsis, including a rapid heartbeat, rapid breathing, or feeling very unwell
\end{itemize}

\section*{Diagnosis and Investigations}
Diagnosis is usually clinical, with imaging used to define the extent of the infection.

\begin{itemize}
    \item \textbf{History and examination:} recent tooth symptoms, examination of the teeth, assessment of the cheek swelling, and measurement of mouth opening
    \item \textbf{Dental X-rays:} periapical and panoramic radiographs to identify the source tooth and any bone involvement
    \item \textbf{CT scan:} to delineate the extent of the abscess and exclude spread to other facial spaces or the orbit
    \item \textbf{Ultrasound:} may help confirm a fluid collection within the cheek
    \item \textbf{Blood tests:} full blood count and inflammatory markers when infection is severe
    \item \textbf{Microbiology:} pus sent for culture and sensitivity when drained, to guide antibiotic choice
\end{itemize}

\section*{Management}
Treatment combines surgical drainage of pus with antibiotics and definitive treatment of the source tooth.

\begin{itemize}
    \item \textbf{Drainage:} the pus is released, usually through an intraoral incision or occasionally through the skin, and a drain may be left in place briefly
    \item \textbf{Antibiotics:} commenced intravenously and continued orally, with cover against both oral aerobes and anaerobes
    \item \textbf{Dental treatment:} root canal therapy or extraction of the infected tooth
    \item \textbf{Pain relief:} paracetamol or anti-inflammatory medicines
    \item \textbf{Hospital care:} some patients require admission for intravenous antibiotics and observation
    \item \textbf{Follow-up:} review to confirm the swelling settles and the dental infection is completely resolved
\end{itemize}

\section*{Complications}
Untreated or inadequately drained, the infection can spread beyond the buccal space with serious consequences.

\begin{itemize}
    \item Spread to other facial spaces, the neck, or the retro-orbital tissues
    \item Cavernous sinus thrombosis, a dangerous infection of a major vein at the base of the brain
    \item The abscess pointing through the skin, with resultant scarring
    \item Sepsis and bloodstream infection
    \item Airway obstruction if infection reaches the floor of the mouth or neck
    \item Loss of the affected tooth
    \item Osteomyelitis of the jaw
\end{itemize}

\section*{Prognosis and Outcomes}
With prompt drainage, appropriate antibiotics, and definitive treatment of the source tooth, the great majority of buccal space abscesses resolve completely within days to a few weeks, usually leaving no lasting problem other than possible loss of the affected tooth. Delay in treatment makes complications more likely and recovery longer. Recurrence is unusual when the underlying dental disease is properly managed.

\section*{Prevention and Self-Care}
The best way to prevent a buccal space abscess is to prevent and treat dental decay early.

\begin{itemize}
    \item See a dentist regularly for check-ups, and seek care promptly for toothache
    \item Brush twice daily with fluoride toothpaste and clean between the teeth
    \item Limit sugary foods and drinks
    \item Do not delay treatment of a decayed, painful, or broken tooth
    \item Follow dental advice after root canal treatment or extraction
    \item Manage conditions such as diabetes carefully, as they raise the risk of infection
    \item Seek early review if any swelling develops after dental treatment
\end{itemize}

\section*{Sources and Further Reading}
The following organisations provide reliable information on dental infections, facial space abscesses, and emergency dental care.

\begin{itemize}
    \item NHS. \textit{Abscess (dental): symptoms, treatment, and when to seek urgent care.} https://www.nhs.uk/conditions/abscess/
    \item Mayo Clinic. \textit{Tooth abscess: causes, treatment, and prevention.} https://www.mayoclinic.org/
    \item MedlinePlus. \textit{Tooth abscess and dental emergencies health information.} https://medlineplus.gov/
    \item MSD Manual. \textit{Abscesses of the orofacial region and dental infection complications.} https://www.msdmanuals.com/
    \item Australian Dental Association. \textit{Dental emergencies and oral health information.} https://www.teeth.org.au/
    \item National Institute for Health and Care Excellence. \textit{Antimicrobial prescribing for dental infections.} https://www.nice.org.uk/
\end{itemize}
